% !TEX root = main_项目后期研究方案.tex
% 由 项目后期研究方案.md 转换；保持严谨求实的工科语言风格与规范的技术路线排版。

\begin{center}
{\zihao{3}\bfseries GearPro 齿轮外观缺陷检测系统：项目后期研究方案\par}
\end{center}

\vspace{0.6em}
\noindent\textbf{报告阶段}：后期研究规划与实施方案 \hfill
\textbf{报告成员}：软件与算法开发组（大三学生） \hfill
\textbf{指导教师}：项目指导教师 \par
\vspace{0.8em}

\textbf{项目阶段进展与后期方案简要概括}：在前期研究中，团队完成了 GearPro 第一代实验台架搭建、50 个真实齿轮的多角度数据建库以及算法离线验证，打通了“输送带送料 $\rightarrow$ 双相机采集 $\rightarrow$ Orin NX 本地推理 $\rightarrow$ Web 监控 $\rightarrow$ 串口通信 $\rightarrow$ STM32 步进电机旋转分拣”的全物理闭环，并在 150 张离线测试集上实现了 88.73\% 的缺陷检出率。本方案针对前期总结出的硬件打光单一、工件 ID 未隔离、单图误报偏高、双相机软件联合判定待完成以及边缘端多模型深度部署探索不足等具体短板，系统制定了项目后期的深入研究方案。后期研究将围绕“光度立体多光源频闪硬件研制、新工件独立盲测集构建、连续多帧时序滑动窗口去重、双相机工件级软件取或、Jetson Orin NX 深度 INT8 量化部署以及人机协同复核规程落地”五大模块循序展开，力争交付一套软硬件闭环稳定、整件误报可控、满足流水线节拍的齿轮外观质检实用原型。

\section*{一、 后期研究总体目标与技术路线}

\subsection*{1. 总体研发目标}
在前期完成的第一代实物台架搭建、6 轮算法实验演进以及 150 张离线测试集评测（检出率 88.73\%、误报率 22.78\%）的基础上，后期的核心目标是：\textbf{推动检测系统从当前的“实验室离线单图验证”平稳过渡到“流水线工件级连续稳定运行”}。

具体而言，后期研究将针对前期暴露出的实际局限，力求达成以下具体目标：
\begin{enumerate}
    \item \textbf{成像硬件升级}：研制基于光度立体视觉的多光源受控频闪成像系统，从物理源头抑制金属反光干扰；
    \item \textbf{数据独立性验证}：采购全新批次齿轮建立实体工件 ID 数据库，提供客观可靠的独立泛化证据；
    \item \textbf{判定逻辑完善}：在软件中实现双相机工件级联合判定，并引入连续多帧时序滑动窗口去重算法，将整件误报率控制在较低水平；
    \item \textbf{边缘端深度部署}：将核心缺陷识别多模型在 Jetson Orin NX 上实施深度 TensorRT 量化与多线程调度优化，实测长时温升稳定性与连续生产节拍；
    \item \textbf{作业规程规范}：完善人机协同复核操作流程，形成闭环质量控制规范。
\end{enumerate}

\subsection*{2. 总体技术路线规划}
后期研究的技术路线紧密围绕“\textbf{多光源成像改造 $\rightarrow$ 新工件盲测建库 $\rightarrow$ 多帧时序去重与双路判定 $\rightarrow$ 边缘端 INT8 量化 $\rightarrow$ 闭环联调与复核规范}”五个有机环节展开，技术实施路线演进关系如图所示：

\begin{center}
\begin{minipage}{0.96\linewidth}
\begin{lstlisting}[basicstyle=\zihao{-5}\ttfamily, frame=single, rulecolor=\color{gray!40}, backgroundcolor=\color{gray!4}, captionpos=b, xleftmargin=0.5em, xrightmargin=0.5em]
                 【后期研究总体实施路线演进图】
+-------------------------------------------------------------+
| 1. 光度立体成像系统研制 (四方向可控频闪 + 纹理/梯度/曲率解耦) |
+------------------------------+------------------------------+
                               |
                               v
+-------------------------------------------------------------+
| 2. 独立工件盲测集构建 (采购全新批次齿轮 + 物理工件 ID 严格隔离) |
+------------------------------+------------------------------+
                               |
                               v
+-------------------------------------------------------------+
| 3. 多帧时序去重与双相机软件判定 (滑动窗口去重 + 双路画面取或)   |
+------------------------------+------------------------------+
                               |
                               v
+-------------------------------------------------------------+
| 4. Jetson Orin NX 深度量化部署 (多模型 TensorRT/INT8 + 异步调度)|
+------------------------------+------------------------------+
                               |
                               v
+-------------------------------------------------------------+
| 5. 整机长时运转压测与复核 SOP (连续运转温升 + REJECT复核/抽检)|
+-------------------------------------------------------------+
\end{lstlisting}
\end{minipage}
\end{center}

\section*{二、 重点研究内容与技术实施方案}

\subsection*{模块一：基于光度立体视觉（Photometric Stereo）的多光源成像系统研制}
\begin{itemize}
    \item \textbf{研究背景与针对问题}：
    前期实验表明，金属齿轮表面的强烈镜面反射亮斑与暗影是引起微细划痕漏检和正常刀纹误报的主要物理根源。第一代台架仅采用了稳态白色漫射光源，尚无法从物理成像源头将表面真实几何形貌缺陷与单纯的光照亮斑区分开来。
    \item \textbf{文献依据与具体方案}：
    通过研读 Saiz 等人与 Wu 等人关于高反光金属工件外观检测的研究，我们了解到光度立体视觉是分离几何凹凸与反射光斑的成熟有效手段。其实施方案规划如下：
    \begin{enumerate}
        \item \textbf{硬件模组设计与装配}：在工业相机镜头周围设计并装配四方向可控的高亮度 LED 环形光源阵列，由 STM32 单片机输出微秒级 PWM 脉冲进行精确的高速频闪同步控制；
        \item \textbf{时序拍摄逻辑}：当输送带将齿轮送至透明亚克力载台并处于短暂平稳状态时，触发单片机控制四方向光源快速交替点亮，工业相机在极短时间内连续抓拍 4 帧不同打光角度下的图像；
        \item \textbf{算法解耦与多通道融合}：算法端基于光度立体法向量方程，解耦计算出原始灰度纹理图（Texture）、反映局部边缘剧烈变化的图像梯度幅值范围图（Range）以及表面平均与高斯曲率图（Curvature）；随后将这三通道特征合成为单张 RGB 伪彩色图像输入下游网络。由于纯反射亮斑在梯度和曲率图上不会产生几何深度突变，而真实划痕与磕碰缺损具有明显的物理凹陷与边缘梯度，该方法能够在物理源头有效剔除反光虚警。
    \end{enumerate}
\end{itemize}

\subsection*{模块二：新批次齿轮物理工件 ID 严格隔离与独立盲测集构建}
\begin{itemize}
    \item \textbf{研究背景与针对问题}：
    前期自建的 50 个真实齿轮数据集是以图像编号进行管理的。尽管前期采用了 64 位感知哈希（dHash）进行了近重复组聚类防泄漏，但尚未在数据库底层建立严格的实体工件 ID 溯源。目前的 88.73\% 召回率是在现有测试图片上的表现，面对未见过的全新批次齿轮时的泛化能力仍需补充独立测试证据。
    \item \textbf{具体实施方案}：
    \begin{enumerate}
        \item \textbf{实体工件 ID 数据库重构}：对现有 50 个齿轮的历史照片进行全面复核，建立 \texttt{Gear\_Entity\_ID} $\rightarrow$ \texttt{View\_Images} 的底层实体溯源标签；
        \item \textbf{增购全新批次实物齿轮}：团队计划采购与定制一批完全未参与过此前训练的全新齿轮（包括 10 个正常齿轮、10 个划痕齿轮、10 个缺损齿轮，共 30 个全新实物）；
        \item \textbf{标准盲测实验设计}：在不微调、不重新训练现有模型参数的前提下，将新采购齿轮送入台架进行图像采集与全自动测试，严格统计基于物理实体工件维度的真实检出率与误报率，为算法泛化性能提供客观证据。
    \end{enumerate}
\end{itemize}

\subsection*{模块三：旋转多帧时序滑动窗口去重与双相机工件级判定算法研发}
\begin{itemize}
    \item \textbf{研究背景与针对问题}：
    前期测试表明单张图像的误报率为 22.78\%，且台架虽然在物理上安装了上下两台工业相机，但在现阶段软件流程中主要还是基于单路相机独立图像进行测试，双相机各自检测后“取或”的工件级软件判定逻辑尚未闭环。
    \item \textbf{具体实施方案}：
    \begin{enumerate}
        \item \textbf{多帧时序滑动窗口去重算法}：借鉴 Allam 等人在汽车齿轮检测中的连续多帧时序判定规则，在上位机软件中设计跨帧时序过滤模块。齿轮在旋转分度过程中，单帧检出的疑似缺陷仅标记为候选；系统在时序角度上维护滑动窗口，\textbf{只有当特定齿廓区域在连续 $N$ 帧（如连续 2$\sim$3 帧）中均被检测出疑似缺陷时，才真正触发工件级的不合格报警}，以此过滤单帧偶然反光带来的虚警；
        \item \textbf{双相机软件联合判定闭环}：完善上位机多线程相机驱动，实现上下两台 2560 $\times$ 1440 工业相机的同步触发采图；上相机负责齿轮端面与齿顶轮廓，下相机负责齿面渐开线与齿根；两路画面分别经过 ROI 提取与缺陷识别后，在软件层实施“工件级逻辑或（Workpiece-level OR）”合并判定，彻底消除单视角遮挡死角。
    \end{enumerate}
\end{itemize}

\subsection*{模块四：Jetson Orin NX 嵌入式边缘端深度量化部署与节拍压测}
\begin{itemize}
    \item \textbf{研究背景与针对问题}：
    前期仅对齿轮定位模型完成了 TensorRT FP16 的基础加速转换（单环节 P95 耗时由 65.0 ms 降至 26.7 ms），核心缺陷识别的多模型融合系统（含 P2-YOLO 与分类网络）仍运行在 PyTorch 框架下，整体部署端的技术探索较少，尚未测量整套多模型流水线在边缘端的实际推理延迟、显存占用与长时温升。
    \item \textbf{具体实施方案}：
    \begin{enumerate}
        \item \textbf{多模型深度 TensorRT 量化编译}：将划痕子系统与缺齿子系统的全部 YOLO 检测网络与分类网络导出为 ONNX 格式，在 Jetson Orin NX 目标板上利用 TensorRT 工具链进行 FP16 半精度（及校准集 INT8）量化编译，充分调用 Orin NX 的 GPU 与硬件加速核心；
        \item \textbf{异步多线程流水线优化}：设计“双相机 SDK 视频流抓取 $\rightarrow$ 定位 ROI 裁剪 $\rightarrow$ 缺陷多模型并发推理 $\rightarrow$ 时序滑动窗口判定 $\rightarrow$ 串口指令下发”的异步多线程并发架构，消除 I/O 阻塞；
        \item \textbf{长时连续运转与温升压测}：在台架上开展 4 小时以上的连续不间断运转试验，系统记录 Orin NX 的 GPU 利用率、显存峰值、芯片运行温度及功耗曲线，确保设备在长时间工业作业下不发生热降频或内存泄漏，使单件全流程处理时间稳定匹配产线输送节拍。
    \end{enumerate}
\end{itemize}

\subsection*{模块五：机电分拣动作闭环联调与人机协同复核规程（SOP）构建}
\begin{itemize}
    \item \textbf{研究背景与针对问题}：
    系统现阶段定位为辅助人工复核的原型工具，为了使检测结果能够稳定指导现场作业，必须打通可靠的机械动作闭环并建立标准化作业规程。
    \item \textbf{具体实施方案}：
    \begin{enumerate}
        \item \textbf{STM32 分拣动作时序优化}：完善上位机与 STM32 的硬件串口（UART，115200 波特率）握手通信。当软件判定为 \texttt{02}（REJECT）时，下位机根据光电开关位置与旋转编码器脉冲，精准控制步进电机平滑旋转至疑似缺陷工位，并驱动执行机构平稳接料；
        \item \textbf{Web 端人机交互体验优化}：在 Web 监控界面中增加“缺陷局部放大”、“历史报警图像回溯”与“判定置信度展示”功能，便于质检员在现场直观复核；
        \item \textbf{人机协同复核规程（SOP）构建}：
        \begin{itemize}
            \item \textbf{REJECT 疑似件复核}：质检工人根据屏幕红框提示，在 1$\sim$2 秒内对实物工件进行快速目视确认，属实则作为不合格品流转，若为误报则一键确认放行；
            \item \textbf{PASS 件定期抽检机制}：设定固定比例（如 5\%）的定期人工抽检流程，监控系统是否存在未知漏检，形成全流程闭环质量保障。
        \end{itemize}
    \end{enumerate}
\end{itemize}

\section*{三、 预期交付成果与量化考核指标}

\subsection*{1. 预期交付的成果清单}
\begin{enumerate}
    \item \textbf{数据集成果}：包含严密实体工件 ID 标签的自建数据集，以及一份完全独立的《30 个全新实物齿轮独立盲测数据集》；
    \item \textbf{硬件成果}：一套升级了四方向受控频闪光度立体照明、上下双相机视野覆盖、透明亚克力载台与步进电机平滑分拣的第一代整机改进实物样机；
    \item \textbf{软件与算法成果}：基于 TensorRT 全面加速的嵌入式边缘端检测软件（含光度立体解耦模块、多帧时序滑动窗口去重模块及 Web 监控系统）；
    \item \textbf{工程技术文档}：包含《新齿轮工件级独立盲测技术报告》、《Jetson Orin NX 边缘部署性能与长时稳定性测试报告》及《齿轮外观质检人机协同复核 SOP 手册》。
\end{enumerate}

\subsection*{2. 量化技术考核指标}

\begin{table}[!htbp]
\centering
\caption*{\textbf{表 1　GearPro 齿轮外观缺陷检测系统后期研究量化技术指标}}
\small
\renewcommand{\arraystretch}{1.25}
\setlength{\tabcolsep}{6pt}
\begin{tabularx}{\linewidth}{@{}>{\bfseries\raggedright\arraybackslash}p{3.5cm}>{\centering\arraybackslash}p{3.2cm}>{\centering\arraybackslash}p{3.2cm}>{\raggedright\arraybackslash}X@{}}
\toprule
考核维度 & 当前前期状态 & 后期预期指标 & 验证与测试方法 \\
\midrule
独立工件盲测召回率 & 88.73\%（仅限单图） & \textbf{$\ge$ 90.0\%（工件级）} & 在 30 个全新未见齿轮盲测集上实测统计 \\
整件综合误报率 (FPR) & 22.78\%（单图误报） & \textbf{$\le$ 8.0\%（工件级）} & 结合多帧时序滑动窗口规则后实测统计 \\
单件全流程检测耗时 & 定位 26.7 ms（缺陷未测） & \textbf{$\le$ 1.5 秒 / 齿轮} & Orin NX 端端全链路实测（含抓拍、推理与分拣） \\
长时连续运行稳定性 & 仅短时间功能测试 & \textbf{连续运转 $\ge$ 4 小时} & 台架连续运行无崩溃、无内存泄漏、无热降频 \\
机械分拣执行响应率 & 初步物理闭环 & \textbf{100\% 正确触发} & 软件 REJECT 指令与步进电机分拣动作一一对应 \\
\bottomrule
\end{tabularx}
\end{table}

\section*{四、 实施进度计划与阶段里程碑}

为确保后期研究工作扎实有序推进，团队制定了为期 12 周的详细进度计划：

\begin{center}
\begin{minipage}{0.96\linewidth}
\begin{lstlisting}[basicstyle=\zihao{-5}\ttfamily, frame=single, rulecolor=\color{gray!40}, backgroundcolor=\color{gray!4}, captionpos=b, xleftmargin=0.5em, xrightmargin=0.5em]
                      【后期研究实施进度甘特图】
 周次  1-2周   3-4周   5-6周   7-8周   9-10周  11-12周
  +---[模组研制]---+                                (光度立体硬件与频闪调试)
          +---[数据构建]---+                        (采购新齿轮与盲测集标定)
                  +---[算法研发]---+                (时序滑动窗口与双路判定)
                          +---[边缘优化]---+        (Orin NX TensorRT/INT8 量化)
                                  +---[整机压测]---+ (4小时连续运转与温升实测)
                                          +---[总结结题] (撰写盲测报告与复核 SOP)
\end{lstlisting}
\end{minipage}
\end{center}

\begin{enumerate}
    \item \textbf{第 1 $\sim$ 2 周（光度立体模组研制与打光测试）}：设计并装配四方向可控 LED 频闪光源模组，编写 STM32 微秒级脉冲同步控制程序，完成基础打光与法向量解耦测试；
    \item \textbf{第 3 $\sim$ 4 周（新齿轮采购与独立盲测集建库）}：采购 30 个全新批次齿轮，规范建立物理工件 ID 标签库，采集多角度光度立体图像并完成标注；
    \item \textbf{第 5 $\sim$ 6 周（多帧时序去重与双相机软件判定开发）}：编写上位机双相机多线程同步采图模块，在软件中固化时序滑动窗口去重算法与工件级取或判定逻辑；
    \item \textbf{第 7 $\sim$ 8 周（Jetson Orin NX 深度 TensorRT 量化部署）}：将缺陷识别多模型全部导出并完成 TensorRT FP16/INT8 引擎编译，优化多线程异步调度流水线；
    \item \textbf{第 9 $\sim$ 10 周（整机长时连续运转与生产节拍压测）}：在台架上开展连续 4 小时以上的长时运转试验，记录系统处理节拍、Orin NX 温升曲线及分拣机构动作响应率；
    \item \textbf{第 11 $\sim$ 12 周（人机协同复核规程制定与结题总结）}：在车间现场模拟质检员复核流程，完善 Web 控制台交互界面，制定《复核 SOP 手册》，汇总形成全套结题技术材料。
\end{enumerate}

\section*{五、 实训总结与后续展望}

前期扎实的研究与实训验证了 GearPro 系统在打光设计、分缺陷建模、温度校准及第一代整机实物闭环上的可行性。面对前期总结出的客观短板，本后期研究方案不回避问题、不夸大成果，始终立足于大三学生的工程实践实际。通过系统推进光度立体多光源成像研制、新工件独立盲测、时序滑动窗口去重及边缘端深度量化部署，团队有信心在指导老师的带领下，扎实攻克工程细节，力争在结题前将系统打磨为兼具学术支撑与工业实用价值的质检样机成果。
